\documentclass[11pt,a4paper]{article}

\usepackage[margin=1in]{geometry}
\usepackage{amsmath, amssymb, amsthm}
\usepackage{microtype}
\usepackage[hidelinks]{hyperref}
\usepackage{parskip}

\title{Three Insights Toward a Theory of \\ Compositional Generalization in Learned Systems}
\author{}
\date{}

\begin{document}
\maketitle

\begin{abstract}
\noindent
This note summarizes three connected insights toward a unified account of representation,
generalization, and inductive bias in learned systems. The first frames data as a
\emph{holonarchic} structure whose elements have a \emph{dispositional} character,
operationalized through conditional Kolmogorov complexity expressed as a function of
representational radius. The second derives from this a procedure for \emph{mapping the
frontier of known observations} and \emph{estimating the location of novel ones} via
threshold-based reconstruction and recursive association decomposition. The third is a
phenomenological critique of the smoothness priors that dominate current regularization
practice, arguing that they may reflect anthropic biases about cause and effect rather
than properties of the data manifold itself.
\end{abstract}

\section{Holonarchy and Dispositional K-Complexity}

A data distribution $D$ is treated not as a flat set of points but as a
\emph{holonarchic} structure: a hierarchy in which each element is simultaneously a whole
(with internal compositional structure) and a part (entering as a component of larger
wholes). Under this view, an observation $x$ is not adequately characterized by its
features alone --- it is characterized by its \emph{dispositions}: the set of
compositional roles it can play across levels of the hierarchy.

We make this precise through conditional Kolmogorov complexity. Define
\[
  K(x \mid D, r) \;\;=\;\; \text{minimal description length of } x \text{ given the
  subset of } D \text{ within representational radius } r \text{ of } x.
\]
The function $r \mapsto K(x \mid D, r)$ is the \emph{dispositional curve} of $x$. Its
shape encodes how the explanatory burden of $x$ shifts as the conditioning neighbourhood
grows. Three regimes are notable:

\begin{itemize}
  \item For small $r$, $K(x \mid D, r)$ is large: the local neighbourhood is too small
        to compress $x$ effectively.
  \item For intermediate $r$, $K(x \mid D, r)$ drops sharply at the radius at which
        $x$'s compositional structure becomes specifiable in terms of nearby holons.
  \item For large $r$, $K(x \mid D, r)$ asymptotes to a floor set by the prior
        complexity of the full data distribution.
\end{itemize}

The position of the inflection encodes the level of the holonarchy at which $x$ is most
naturally described. The disposition of $x$ is the integral of associations that become
available as $r$ grows past this inflection.

\section{Frontier Mapping and Generative Estimation via Two Thresholds}

The dispositional view suggests an operational procedure for characterising any
observation: identify two critical thresholds in a noise- or perturbation-conditioned
reconstruction process.

Let $\tilde{x}_t$ be a perturbed version of $x$ at noise level $t \in [0, T]$, and let
$\hat{x}_t$ denote the reconstruction by a model conditioned on $D$. Define:

\begin{itemize}
  \item $T_1(x)$: the smallest $t$ at which $x$-specific reconstruction fails ---
        i.e.\ the noise level at which the model can no longer recover the specific
        instance $x$ and must hedge across nearby holons. This marks the
        \emph{instance-to-class transition}.
  \item $T_2(x)$: the smallest $t \geq T_1(x)$ at which the reconstruction collapses to
        the unconditional prior of $D$. This marks the \emph{class-to-prior transition}.
\end{itemize}

From these we derive two quantities of interest:

\[
  \text{Novelty}(x) \;=\; T_2(x) - T_1(x), \qquad
  \text{Associativity}(x) \;=\; T_2(x).
\]

Novelty measures the width of the regime in which $x$ retains class-level identity
without being fully specified --- the size of the compositional ``neighbourhood'' it
governs. Associativity measures how deeply $x$ is anchored in the data manifold before
the reconstruction loses all conditional information.

These two scalars define a coordinate system on the data: every observation maps to a
point in the $(T_1, T_2)$ plane. The frontier of known observations is the boundary in
this plane beyond which no training data lies, and the estimated location of a new
observation can be obtained by:

\begin{enumerate}
  \item Computing its $(T_1, T_2)$ coordinates from the reconstruction process.
  \item Identifying the nearest known holonarchic region in this plane.
  \item \textbf{Recursively decomposing} the associations active at $T_2$ ---
        i.e.\ examining which subsets of $D$ remain influential at the class-to-prior
        boundary, and applying the same threshold procedure to each subset.
\end{enumerate}

This recursion is the \emph{association deconstruction} step: by repeatedly applying the
two-threshold scoring to progressively smaller and more specialized subsets, we descend
the holonarchic hierarchy and obtain a multi-resolution characterization of $x$ in
terms of compositional dependencies on $D$. The scoring is not on $x$ in isolation but
on \emph{which hierarchical decompositions remain consistent} with $x$ across levels.
This is what should be scored: not features, not reconstructions, but the structure of
the hierarchy of associations that survive perturbation.

\section{A Phenomenological Critique of Smoothness Priors}

The dominant regularization practices in deep learning --- dropout, $L_2$ weight decay,
input noise --- share a common implicit assumption: that the function spaces we wish to
learn are \emph{smooth}, in the sense that small perturbations to inputs or weights
correspond to small changes in output. This assumption is so deeply embedded that it is
often not recognized as an assumption at all.

We argue that this commitment is not metaphysically necessary but reflects the
phenomenological structure of human cognition. The way humans perceive and interact with
reality --- the categories of cause and effect, the expectation of continuous response,
the framing of events as smoothly unfolding in time --- biases scientific and
mathematical formalization toward smooth solution spaces. The smoothness prior is in
this sense an \emph{anthropic projection}: it describes the kinds of regularities that
human cognition is equipped to perceive and act upon, not necessarily the structure of
the data manifold itself.

This critique opens a question about what alternative priors are available. Sparsity
priors (in the Kolmogorov/Occam sense) commit to a different assumption: that good
explanations admit compact compositional descriptions, but make no commitment to
continuity. Compositional priors --- such as the holonarchic structure described in
Section~1 --- commit to a third position: that good explanations decompose into typed
hierarchical associations whose admissible compositions are constrained but not
necessarily smooth. The threshold-based scoring of Section~2 is compatible with
compositional priors without requiring smoothness.

The practical consequence is that regularization should not be chosen by default but
by commitment to a theory of what regularities the system is meant to discover. If the
target is compositional generalization across a holonarchic structure, the regularizer
should constrain admissible compositions, not penalize discontinuity.

\section*{Synthesis}

The three insights form a single argument. Holonarchic structure and dispositional
K-complexity (Section~1) provide an ontology of what data \emph{is}; threshold-based
reconstruction and recursive association decomposition (Section~2) provide an
epistemology of how to characterize and place observations within that ontology; and the
phenomenological critique (Section~3) clears the ground of inherited regularization
commitments that may be incompatible with both. Together they sketch a research program
in which the choice of inductive bias is itself a theoretical commitment about the
nature of compositional generalization, rather than a hyperparameter to be tuned.

\section{Paper Outline (Placeholder)}

This appendix sketches the planned structure of the full paper and the experimental
program that operationalizes the three insights above. All architectural variants will
be reported at a common parameter budget (target: $\approx 50$M) so that comparisons
isolate the effect of inductive bias rather than capacity.

\subsection*{Planned Sections}

\begin{enumerate}
  \item \textbf{Introduction.} Motivate compositional generalization as the central
        problem and position the paper's contribution as an examination of inductive
        bias rather than scaling.
  \item \textbf{Theoretical motivation.} The three insights of this document, expanded
        with formal definitions of $K(x \mid D, r)$, the two-threshold reconstruction
        scoring, and the phenomenological critique.
  \item \textbf{Methods.} The architectural progression and training schemes:
        \begin{enumerate}
          \item Flat baseline transformer (AdamW, no curriculum).
          \item Iterative depth-growth training (iter scheme).
          \item Iterative + data curriculum (difficulty-ranked exposure).
          \item Iterative + curriculum + decaying intra-layer loops with horizon-1
                BPTT and PCGrad over loop losses (\emph{loopy iter}).
          \item Sparsity-prior post-training: iterative magnitude- or activation-based
                sparsification of the trained dense model, schedule-tuned until the
                train--validation gap reaches a target.
          \item Holonarchic-prior architecture (\emph{diamond}): branch
                specialization with direction-in-representation-space assignment and
                radius-as-depth structure, instantiating the holonarchic ontology of
                Section~1.
        \end{enumerate}
  \item \textbf{Experiments.} Common-budget comparison across all variants on byte-level
        \texttt{enwik8}, reported as bits-per-byte (bpb) on the held-out set, with
        train--val gap reported alongside.
  \item \textbf{Catastrophic-forgetting study.} A targeted experiment to demonstrate
        the practical value of the holonarchic prior: post-training fine-tuning of a
        single semantic branch (e.g.\ a mathematics branch) of the diamond model and
        measurement of performance preservation on unrelated branches, compared
        against the same fine-tuning applied to a flat baseline of equal parameter
        count. The prediction is that the diamond confines fine-tuning updates to the
        targeted branch and the shared prefix layers, yielding substantially lower
        degradation on distant capabilities.
  \item \textbf{Discussion.} What the comparison reveals about the relative merits of
        sparsity priors versus compositional/holonarchic priors, and what the
        catastrophic-forgetting result implies for modular post-training.
  \item \textbf{Conclusion.}
\end{enumerate}

\subsection*{Model Variants}

\begin{center}
\begin{tabular}{lll}
\hline
\textbf{Variant} & \textbf{Inductive bias} & \textbf{Status} \\
\hline
Flat baseline (AdamW)           & none (control)                 & trained (small) \\
Iter                            & depth curriculum               & trained \\
Iter + data curriculum          & depth + difficulty curriculum  & trained \\
Iter + curriculum + loops       & + intra-layer recurrence       & implemented \\
Sparsity-pruned (TurboSparse)   & sparsity prior                 & planned \\
Diamond                         & holonarchic prior              & planned \\
Diamond + branch fine-tune      & forgetting resistance test     & planned \\
\hline
\end{tabular}
\end{center}

\subsection*{Key Empirical Claims to Be Tested}

\begin{itemize}
  \item Iterative depth-growth with data curriculum converges faster per training step
        than a matched-capacity flat baseline and reaches comparable bpb.
  \item Sparsity-prior post-training reduces the train--val gap without recourse to
        smoothness-based regularizers, providing an empirical test of the
        phenomenological argument of Section~3.
  \item The diamond architecture achieves bpb competitive with flat baselines at equal
        parameter count, demonstrating that the holonarchic inductive bias is not
        prohibitively expensive.
  \item Branch-localized fine-tuning of the diamond preserves capabilities on distant
        branches substantially better than equivalent fine-tuning of a flat model.
        This is the headline result that links the theoretical claims of
        Sections~1--3 to a measurable practical property.
\end{itemize}

\end{document}
